\begin{abstract}
Agents that reward hack often explore their environment before exploiting it: in the 2026 OpenAI--Hugging Face incident, environment probing rose during training of the model that later attacked a production system~\citep{openai2026report}. We ask whether reward-hacking trajectories show more environment probing than legitimate ones, measured \emph{only from the agent's commands}, without access to its reasoning. We define a rule-based taxonomy of probing (reading the verifier, exploring outside the task scope, system and runtime introspection, reaching the outside world) and score each trajectory by its number of probing commands. On all 8{,}641 trajectories of Terminal Wrench~\citep{terminalwrench}, probing separates attack-prompted agents from baseline agents (Claude Opus 4.6: AUC 0.68 after controlling for trajectory length; 248 of 295 tasks), but it does not separate successful hacks from legitimate solves produced under the same attack prompt (AUC 0.48--0.50 across three models and three robustness controls). On EvilGenie~\citep{evilgenie}, where no one asks the agent to cheat, an open 20B model reward hacked spontaneously in 1 of 333 attempts, after reading the test file. Command-level probing is thus an early signal of an agent \emph{looking for} shortcuts, not evidence that it \emph{took} one. \todo{Agregar resultado de la validación humana (kappa) cuando esté.}
\end{abstract}
